\section{Proofs for the ladder theorems}\label{app:packdef}

\begin{proof}[Proof of \Cref{lem:idempotent-iterates}]
For $n=1$ there is nothing to prove.
If $e^{\circ n}=e$, then $e^{\circ(n+1)}=e^{\circ n}\circ e=e\circ e=e$.
\end{proof}

\begin{proof}[Proof of \Cref{lem:definability-bijection}]
\emph{Well defined and onto.}
For $A\subseteq\im(f)$, the predicate $\1_A\circ f$ is $f$-definable, with $g=\1_A$.
Conversely, let $h=g\circ f$ be $f$-definable and put $A=\{y\in\im(f):g(y)=1\}$.
For every $x$, $f(x)\in\im(f)$, so $h(x)=1$ if and only if $f(x)\in A$; thus $h=\1_A\circ f$.

\emph{Order isomorphism.}
If $A\subseteq B$ then $\1_A\le\1_B$, hence $\1_A\circ f\le\1_B\circ f$.
Conversely, suppose $\1_A\circ f\le\1_B\circ f$ and let $y\in A$.
Since $y\in\im(f)$ there is $x$ with $f(x)=y$, and $1=\1_A(f(x))\le\1_B(f(x))$, so $y\in B$.
Thus the map reflects the order, and in particular it is injective: $\1_A\circ f=\1_B\circ f$ implies $A\subseteq B$ and $B\subseteq A$.
\end{proof}
